\scanpage{16}
\mspage{14}
à $X'$. Lorsque $k=k'$, $S=S'$, et que les morphismes envisagés sont les puissances $p$.èmes, on est donc conduit à chercher à \deleted{prolonger} relever ce morphisme à $X$ de façon compatible avec le morphisme $\operatorname{Specf}(W)\longrightarrow\operatorname{Specf}(W)$ déduit de $f_W$, ou ce qui revient au même, de relever le morphisme canonique $S\longrightarrow S^{(p/k)}$ en un morphisme de $R$-schémas formels $X\longrightarrow X\times_R(R,f_R)$. C'est en tous cas toujours possible si $S$ est affine, cas auquel on peut se ramener. Ayant donc ainsi un morphisme
\[
f_X:X\longrightarrow X,
\]
compatible avec $f_R$ sur $R=\operatorname{Specf}(W)$, et induisant $f_S$ sur $S$, le foncteur $\underline{M}\rightsquigarrow\underline{M}^{(p)}$ s'explicite comme l'image inverse ordinaire de Modules stratifiés sur $X$ relativement à $R$, resp. (dans le cas de isocristaux) de Modules stratifiés sur l'espace rigide-analytique $X_K$, pour le ``morphisme'' $X_K\longrightarrow X_K$ d'espaces rigide-analytiques (relatif à $f_K$ sur le corps de base !) induit par $f_X$. Bien que $f_X$ et par suite $f_{X_K}$ \unsure{soit} loin d'être unique, le foncteur envisagé qu'il définit ne dépend pas, essentiellement, des choix faits.

1.15. \uline{Bicristaux et biisocristaux sur une base quelconque}. Ne supposons plus nécessairement $S$ de caractéristique déterminée $p$. \deleted{Nous} Pour chaque nombre premier $p$, soit $S_p$ la fibre $V(p.1_{\underline{O}_S})$ de $S$ sur le point $p\underline{Z}$ de $\operatorname{Spec}(\underline{Z})$. Un bicristal de poids $i$ sur $S$ est par définition la donnée d'un cristal de modules $\underline{M}$ sur $S$, et pour chaque $p$ d'une structure de bi-cristal de poids $i$ sur la restriction $\underline{M}_p$ de $\underline{M}$ à $S_p$, définie par la donnée de $F_p,V_p$. On définit de même la notion de bi-isocristal sur $S$, étant entendu qu'un isocristal est un objet de la catégorie des isocristaux sur $S$ (comme de juste), définie à partir de la catégorie des cristaux de modules en ``localisant'' par rapport à des entiers $n>0$ arbitraires, i.e. en tensorisant les $\operatorname{Hom}$ sur $\underline{Z}$ par $\underline{Q}$. - Lorsque $S$ est de caractéristique nulle, le supplément de structure impliqué par ``bi'' est évidemment vide, tandis que si $S$ est de caractéristique $p$, on retrouve les notions de 1.10. Lorsque enfin $S$ est de type fini sur $\underline{Z}$ et domine $\operatorname{Spec}(\underline{Z})$, alors la notion envisagée est d'une nature essentiellement arithmétique, les différents $F_p$ jouant le rôle d'homomorphismes de frobenius, comme de bien entendu; dans le cas du poids 1, en
\margin{15}
